Lors d'une répétition générale dans le département de musique, le deuxième mode de vibration d'une corde de contrebasse de longueur $L_c = \CordeLongueur\text{ cm}$ fixée aux deux extrémités et possédant une masse linéique de $\CordeMasseLin\text{ kg/m}$ résonne exactement à la même fréquence que le troisième mode d'un tuyau d'orgue fermé à une extrémité et ouvert à l'autre de longueur $L = \TuyauLongueur\text{ cm}$. La vitesse du son dans la pièce est évaluée à $\CordeVitesseSon\text{ m/s}$.

\begin{enumerate}[label=\alph*)]
    \item Tracez l'allure générale du deuxième mode de la corde et du troisième mode du tuyau ouvert-fermé.
    \item Calculez la tension appliquée sur la corde de contrebasse.
    \item La corde vibre (toujours dans son $2^\text{e}$ mode) et émet des ondes sonores dans l'air ambiant. Quelles sont la fréquence et la longueur d'onde de ces ondes sonores perçues dans la pièce ?
\end{enumerate}